% Study sheet: Volume III, Chapter 20, Conclusion.
% Numbers: \val{...} from conclusion_studysheet_numbers.tex, written by conclusion_studysheet_numbers.py.
\documentclass{studysheet}
\input{conclusion_studysheet_numbers}
\sheetvolume{3}
\sheetchapter{20}{Conclusion}
\sheetsources{Numbers from the chapter's mlsysim scenarios (the remediation agent's illustrative workload profile); questions marked assumed use new inputs.}

\begin{document}
% ===================================================================== FRONT: LEARN
\frontside
\begin{bigidea}
Components that pass their own tests can still leave a task unfinished or its effects unaccounted for. Reliable
autonomy comes from a runtime that settles, at every boundary, \hl{what was proposed, what was done, and what
evidence shows the task is complete}. The \hl{H·S·A exposures} (horizon, state, authority) set how much of that
runtime a task needs, and \hl{invariant closure}, the thread through every part, closes the book.
\end{bigidea}

\begin{leftcol}
\sheetsec{Key ideas}
\begin{ideas}
\idea{The contract sizes the guards.}{Exposure sets the obligations; for an irreversible actor the gates
  (approval, version check, health window) outweigh caps on spend.}
\idea{A lost reply is unknown, not failed.}{Query status under the same operation identity, never a fresh
  retry; hold dependent changes; unsettled work stays with its owner.}
\idea{Authority stays with the source.}{The source record outranks logs, memory, indexes and context;
  a cache hit proves reuse, not freshness.}
\idea{The remedy follows the attributed cause.}{Replay varying one factor. Fix context, tools or runtime
  first: new weights cannot supply a missing record or check.}
\idea{A safety case scopes its claims.}{Claim, argument, evidence, assumptions. A log shows order;
  cause needs a controlled comparison or fault injection.}
\idea{Assisting is not adopting.}{A model may draft contracts and tests; only the authorized owner adopts
  them. Co-written tests are not independent evidence.}
\end{ideas}

\vspace{10pt}
\sheetsec{One incident: an owner and a record at each step}
{\centering
\begin{tikzpicture}[sheet]
  \foreach \i/\lab/\own/\sty in {
      0/{Task admitted}/{runtime, under the adopted contract}/accbox,
      1/{Evidence selected}/{context manager: source, version, time}/nbox,
      2/{Change approved}/{model proposes; operator approves}/nbox,
      3/{Intent logged, sent}/{action gate: durable intent, owner}/nbox,
      4/{Lost reply settled}/{settlement, for owner: endpoint status}/hotbox,
      5/{Incident accepted}/{acceptor: contract, settled operation, health window}/accbox}{
    \node[\sty, text width=72pt, align=center, minimum height=14pt] (n\i) at (40,{-\i*22}) {\lab};
    \node[small, anchor=west, text=ink] at (88,{-\i*22}) {\own};}
  \foreach \i [evaluate=\i as \j using int(\i+1)] in {0,...,4}{\draw[arr] (n\i.south) -- (n\j.north);}
\end{tikzpicture}\par}

\end{leftcol}%
\begin{rightcol}
\fcard{Invariant closure: horizon compounding}%
  {\Pr(\text{no violating step}) = (1-\PerStepErrorRate)^{\TurnHorizon}}%
  {$\TurnHorizon$ unchecked turns, each violating independently with probability $\PerStepErrorRate$. A run holds
   at least one: $1-(1-\PerStepErrorRate)^{\TurnHorizon}$.}%
  {Bounds enforced below the model, acceptance on evidence above it.}{}
\fcard{Cost per accepted task}%
  {\CostPerAcceptedTask = \frac{\mathbb{E}[\AttemptCost]}{\AcceptanceProbability}, \qquad
   \text{a change pays if } \MechanismAcceptanceProbability > \frac{\MechanismAttemptCost}{\CostPerAcceptedTask}}%
  {Rejected attempts are charged to accepted ones. A change raising attempt cost to $\MechanismAttemptCost$
   must lift acceptance to $\MechanismAcceptanceProbability$ by a greater relative fraction.}%
  {}{}
\fcard{Climbing the intervention ladder}%
  {\BreakEvenTaskCount = \frac{\AdaptationInvestmentCost}{\BaselineTaskCost-\AdaptedTaskCost}}%
  {A heavier rung costs $\AdaptationInvestmentCost$ once and cuts each task's cost from $\BaselineTaskCost$
   to $\AdaptedTaskCost$ at the same acceptance; it pays only past $\BreakEvenTaskCount$ tasks.}{}{}
\fcard{Expected irreversible loss}%
  {\ExpectedIrreversibleLoss = \sum_{\PlanAction} \Pr(\text{bad } \PlanAction)\,\text{consequence}(\PlanAction)\,(1-\ActionRecoverableFraction)}%
  {Over actions $\PlanAction$; $\ActionRecoverableFraction$ is the share of harm caught and reversed (zero with
   no check ahead of the write).}{}{}
\fcard{Verifier precision}%
  {\VerifierPrecision = \frac{\ValidCandidateRate(1-\VerifierFalseRejectionRate)}{\ValidCandidateRate(1-\VerifierFalseRejectionRate)+(1-\ValidCandidateRate)\,\VerifierFalseAcceptanceRate}}%
  {Valid share of the passing set: $\ValidCandidateRate$ valid candidates, $\VerifierFalseRejectionRate$ of
   them missed, $\VerifierFalseAcceptanceRate$ of invalid ones passed. The base rate dominates.}{}{}
\end{rightcol}

% ===================================================================== BACK: PRACTICE
\backside
\begin{answers}
\ans{Finish}{\val{ctwo} $\div$ \val{paccfrac} $=$ \$\val{cefftwo}. To match \$\val{ceff}, acceptance must reach
  \val{ctwo} $\div$ \val{ceff} $\approx$ \val{pneedtwo}\%, which is impossible: on cost per accepted task alone it never
  pays. Only the expected loss it removes beyond existing gates (at most \$\val{loss} per attempt) can justify it (loss card).}
\ans{1}{(a) \val{qpassinvalid} $\div$ (\val{qpassvalid} $+$ \val{qpassinvalid}) $\approx$ \val{qinvalid}\% invalid: the base
  rate dominates. (b) $\VerifierFalseAcceptanceRate$ below \val{qfamax}\%, a check \val{qfafold}$\times$ sharper;
  raising the valid share is the other lever (precision card).}
\ans{2}{Unknown. Keep the intent, query the endpoint's status under the same operation identity, hold
  dependent changes. An empty lookup proves nothing: unresolved, kept by the operations owner (lost-reply idea).}
\columnbreak
\ans{3}{Only that reusing those tokens' computation is valid; authority stays with the source record. Test
  context first: replay with the new configuration supplied; if that repairs it, fix context assembly, not weights
  (remedy idea).}
\ans{4}{(a) No: traffic may have subsided; the endpoint's status record shows it. (b) No: a log gives order,
  not cause; run a controlled comparison or fault injection. (c) The contract, the approval, the settled operation and
  new version, the health window and policy checks; post-window monitoring stays with the owner (safety-case idea).}
\end{answers}

\begin{leftcol}
\sheetsec{Worked example}
\begin{practicebox}\raggedright\boxtitle{Does more diagnosis pay on the incident?}
One attempt: model calls \$\val{cmodel}, telemetry and change tools \$\val{cenv}, health check
\$\val{cassess}, approver review \val{reviewmin} min at \$\val{reviewrate}/h $=$ \$\val{capproval}. Acceptance
\val{pacc}\%. A wrong change: \val{pbad}\% of attempts, \$\val{impact} if it commits.\par\vspace{5pt}
{\centering
\begin{tikzpicture}[sheet]
  \pgfmathsetmacro\sc{200/\val{cattempt}}
  \pgfmathsetmacro\xa{\val{cmodel}*\sc} \pgfmathsetmacro\xb{(\val{cmodel}+\val{cenv}+\val{cassess})*\sc}
  \draw[hotbox] (0,0) rectangle (\xa,12);
  \draw[nbox] (\xa,0) rectangle (\xb,12);
  \draw[accbox] (\xb,0) rectangle (200,12);
  \node[lab] at ({(\xb+200)/2},6) {approval \val{approvalshare}\%};
  \node[note, anchor=west] at (204,6) {\$\val{cattempt}};
\end{tikzpicture}\par}\vspace{4pt}
\begin{steps}
\item Attempt: \val{cmodel} $+$ \val{cenv} $+$ \val{cassess} $+$ \val{capproval} $=$ \res{\$\val{cattempt}}; model share
  \val{cmodel} $\div$ \val{cattempt} $=$ \res{\val{modelshare}\%}.
\item Per accepted task: \val{cattempt} $\div$ \val{paccfrac} $=$ \res{\$\val{ceff}}.
\item Double the model calls: \val{cattempt} $+$ \val{cmodel} $=$ \$\val{cgrown}; pays only if acceptance exceeds
  \val{cgrown} $\div$ \val{ceff} $=$ \res{\val{pneed}\%}.
\item Risk: \val{pbad}\% $\times$ \$\val{impact} $=$ \res{\$\val{loss}} per attempt with no check ahead of the write,
  about \val{lossmult}$\times$ the attempt.
\end{steps}
\takeaway{Tokens are cheap here; the gates, not the spend caps, carry the risk.}
\end{practicebox}

\vspace{10pt}
\sheetsec{Your turn: finish it}
\begin{practicebox}\raggedright\boxtitle{\assumed{} a second reviewer signs each approval}
Same incident; the extra reviewer adds another \$\val{capproval} per attempt.\par\vspace{4pt}
\textbf{Given:} new attempt \val{cattempt} $+$ \val{capproval} $=$ \$\val{ctwo}.\par
\textbf{You:} find the cost per accepted task at \val{pacc}\% acceptance and the acceptance needed to match
\$\val{ceff}. On what other account could the reviewer still be justified?\par\vspace{30pt}
\end{practicebox}
\end{leftcol}%
\begin{rightcol}
\sheetsec{Common mistakes}
\mistake{``A capable enough model removes the need for runtime boundaries.''}%
  {At \val{epsviolate}\% violating steps, \val{pany}\% of \val{horizonturns}-turn runs hold one: about
   \val{violday} of \val{trajday} runs a day.}{}
\mistake{``The check passed, so the completion claim is proven.''}%
  {With \val{pvalid}\% valid candidates, a check that misses \val{pmiss}\% of valid and passes \val{pfa}\% of
   invalid ones leaves \val{invalid}\% of its passing set invalid.}{}
\mistake{``The model explained why, so it may shorten the health window.''}%
  {Only the contract owner adopts it. With roughly independent samples, half as many widen the latency confidence interval
   \val{widen}$\times$ ($\sqrt{2}$): weaker evidence when it matters most.}{}

\vspace{2pt}
\sheetsec{Check yourself}
\begin{qs}
\item \assumed{} the same check (misses \val{pmiss}\%, passes \val{pfa}\% of invalid) on candidates only
  \val{qvalid}\% valid. (a) What share of the passing set is invalid? (b) What false-accept rate keeps it at
  \val{qtarget}\%?
\item The rollback request times out. What may the runtime conclude, what does it do next, and what if the status
  lookup comes back empty?
\item After the rollback settles, the next prompt reuses its cached prefix, which still shows the old
  configuration as active. What does the cache hit establish, which principle is broken, and which remedy
  do you test first?
\item Latency is healthy right after the rollback. Does that show (a) the change applied, (b) it caused recovery?
  (c) What must the acceptance record bind before the owner closes the incident?
\end{qs}
\end{rightcol}
\end{document}
